\documentclass[11pt,a4paper]{ctexart}
\usepackage[margin=2.4cm]{geometry}
\usepackage{amsmath,amssymb}
\usepackage{fancyhdr}
\usepackage{enumitem}
\linespread{1.18}
\setlength{\headheight}{15pt}
\setlength{\parskip}{0pt}

\pagestyle{fancy}
\fancyhf{}
\fancyhead[L]{\small Team 2309229}
\fancyhead[R]{\small \thepage}
\renewcommand{\headrulewidth}{0.4pt}
\fancypagestyle{plain}{\fancyhf{}\renewcommand{\headrulewidth}{0pt}}

\ctexset{
  section/format = \Large\bfseries\raggedright,
  subsection/format = \large\bfseries\raggedright,
  subsubsection/format = \normalsize\bfseries\raggedright
}
\setcounter{secnumdepth}{3}

\begin{document}

%====================== 摘要页 ======================
\thispagestyle{empty}
\begin{center}
{\small 2023 MCM/ICM\quad 选题 A\quad 队伍控制号 2309229\quad 摘要页（Summary Sheet）}\\[18pt]
{\LARGE\bfseries 抗旱勇士：植物群落}\\[8pt]
{\large （The Warriors against Drought: Plant Communities）}\\[10pt]
\end{center}

\begin{center}
{\large\bfseries 摘\quad 要}
\end{center}

\vspace{6pt}
随着全球变暖以来严重干旱的不断加剧，许多植物群落因缺水而消亡，而另有一些群落却能够在干旱之后恢复。为探究这一现象，我们建立了 GIID 模型来描述植物群落对干旱的响应，并分析物种多样性是否能够帮助群落抵御干旱。

在建立全部模型之前，我们选取了五种典型植物物种，并依据其固有生长率与最适土壤湿度将它们划分为三个类别。此外，我们对模型中所使用的参数进行了估计。

关于任务一：为模拟植物生长与种间竞争，我们首先建立了 Lotka--Volterra 竞争模型。随后，考虑每个物种的最适土壤湿度，我们将各物种的环境容量设定为土壤湿度的函数。最后，我们对土壤湿度的变化过程进行建模，以描述植物与环境之间的相互作用。我们考虑了植物吸水、不规则降雨以及其他因素。我们对这些分析加以汇总，设计出 GIID 模型。随后，我们验证了模型的正确性。结果如图 7 所示。

关于任务二：首先，我们进行定性分析，发现五物种群落在干旱中更为稳定。为了量化群落从物种多样性中获得的益处，我们计算了总生物量、总生物量的变异系数以及 S-W 指数，分别用以评价生态系统的活力、稳定性与均匀度。实验结果表明：随着物种数目的增加，群落的总生物量先增大后减小，而稳定性与均匀度则持续上升。此外，群落需要至少两个物种才能从中获益。结果如图 10 至图 13 所示。

关于任务三：为探究物种类型的作用，我们选取了三种典型的物种组合，并对每种组合分别进行模拟。我们发现，在最适土壤湿度上存在中等差异的组合能够提高群落的总生物量与稳定性。而如果各物种最适土壤湿度之间的差异过小或过大，群落将无法获益。

关于任务四：我们调整了干旱的强度与频率，实验结果表明：当干旱变得更强、更频繁时，物种之间的互补效应也随之增强。此时群落生物量随物种数目的增加而增加。当干旱不那么频繁时，物种之间的竞争更强，物种数目的增加反而导致总生物量下降。物种数目的增加则始终增强群落的稳定性。

关于任务五与任务六：我们将生境丧失与环境污染的影响归因于环境容量的下降。因此我们建立了一个阶跃函数来模拟这一过程。我们发现，生境丧失与环境污染将导致群落总生物量的下降，而在这一过程中多物种群落仍将保持较高的稳定性。接下来，我们阐述了确保群落长期存续所需采取的措施。我们还基于上述发现描述了群落对更大范围环境的影响。

最后，我们改变植物竞争强度、吸水速率以及气候模式，以便对这些参数进行敏感性分析。我们发现模型对上述所有参数都是敏感的。我们还分析了模型的优势与不足，并讨论了如何对模型加以改进。

\vspace{6pt}
\noindent\textbf{关键词：}Lotka--Volterra 模型，生物多样性，不规则干旱，互补效应

\newpage
\pagestyle{fancy}

%%====================== 正文 ======================
\section{引言}

\subsection{问题背景}

近年来，由于全球变暖，世界范围内极端天气事件的发生频率急剧上升。干旱是对我们的生活产生巨大影响的极端天气事件之一。下图采集自美国干旱监测网站 Drought.gov，它反映了 2023 年 2 月 8 日至 2023 年 2 月 14 日期间美国的干旱状况。

从图中我们可以看到，美国 34.55\% 的国土面积处于干旱之中，波及 2.529 亿英亩的农作物以及 7430 万人口。在极端干旱期间，我们很容易发现不同的植物群落表现出不同的行为。草地与农田中的植物几乎全部死亡，而多物种群落受到的影响较小，并在干旱之后繁茂生长。基于上述分析，研究植物如何响应干旱、以及植物群落中物种的多样性如何帮助植物抵御干旱，既是必要的，也是紧迫的。

\subsection{问题的重述}

\begin{itemize}
\item 根据题目要求，我们首先应建立一个数学模型，用以预测植物群落在遭遇各种不规则天气周期时随时间的变化。并且要求我们在干旱周期中考虑不同植物物种之间的相互作用。
\item 其次，我们应将模型推广到长期时间尺度与更大范围的环境之中。
\item 第三，我们需要进一步讨论植物物种数目与植物群落所获益处之间的关系，并探究当物种数目增加时会发生什么。此外，我们还应关注物种类型对我们此前所得结果的影响。
\item 我们还被要求研究在频率更高或更低、变化幅度更大的干旱条件下，物种数目是否会对植物总种群产生同样的作用。随后，我们还应把污染与生境缩减等其他因素纳入考虑。
\item 最后，基于上述全部结果，我们应弄清楚为确保群落的存续需要做些什么，以及植物群落对更大范围的环境有何影响。
\end{itemize}

\subsection{文献综述}

为了研究干旱中的植物群落，我们需要了解模拟植物群落的既有方法以及相应的生态学原理。根据文献 [1, 2, 3, 4]，在模拟方面主要有三种方法，而在原理方面主要有两条原理需要我们使用。如下列图所示。

\subsection{我们的工作}

本题要求我们探究植物与环境之间的相互作用，并度量物种多样性的作用。我们主要开展了以下工作。

\begin{itemize}
\item 在对植物生长、种间竞争以及水资源竞争等若干因素进行分析的基础上，建立了 GIID 模型。
\item 探究了物种数目与物种类型的作用。
\item 探究了干旱频率变化、生境缩减以及环境污染等其他因素的作用。
\end{itemize}

为避免冗长的文字描述，我们用流程图来展示我们的工作，如图 3 所示。

\section{模型准备}

\subsection{假设与依据}

为简化所给问题，我们作出以下基本假设：

\textbf{假设 1：}每种植物的吸水速率与其生物量成正比。

\textbf{依据：}根据文献 [7]，当植物之间发生激烈竞争时，一株植物所获得的每一种资源都与它的生物量成正比。

\textbf{假设 2：}影响植物生长的主要因素是土壤湿度。并且不同植物生长所需的最适土壤湿度各不相同。

\textbf{依据：}在现实世界中，温度、湿度、日照等影响植物生长的因素相互耦合。然而我们只关注极端干旱条件下植物群落的生长。因此，为简化模型，我们忽略其他因素，只考虑土壤湿度。当植物在适度的土壤湿度下生长时，它会生长得很快；而在缺水或被水淹时则无法生长。因此存在一个最适土壤湿度。

\textbf{假设 3：}在适宜条件下，所有植物共享一个共同的环境容量。

\textbf{依据：}大多数植物是生产者，并以相同的方式获取能量（均来自光合作用）。在适宜条件下，它们都有机会吸收生境中的全部太阳能。因此它们能够达到相同的最大生物量，这意味着它们共享相同的环境容量。

\textbf{假设 4：}在我们的研究中，不考虑植物入侵与气候模式突变等偶发事件的影响。

\textbf{依据：}尽管植物生长在自然环境中，但为简化模型，我们假设植物群落处于一个相对封闭的环境之中。因此既不存在入侵物种，也不存在气候模式的大幅快速变化。

\textbf{假设 5：}在植物群落中，唯一的水分来源是降水。

\textbf{依据：}群落获取水分的途径可能多种多样，例如露水、雾以及降水。由于露水与雾所携带的水量很少，我们忽略它们的作用，只考虑降水。

\subsection{符号说明}

主要符号如表 1 所示。

\begin{itemize}
\item $r_i$，单位为 $\mathrm{day}^{-1}$：第 $i$ 种植物的固有生长率；
\item $x_i$，单位为 $\mathrm{kg}$：第 $i$ 种植物的生物量；
\item $Env_i$，单位为 $\mathrm{kg}$：第 $i$ 种植物的环境容量；
\item $W$，无量纲：土壤湿度；
\item $bestW_i$，无量纲：第 $i$ 种植物所需的最适土壤湿度；
\item $\alpha$，无量纲：竞争强度；
\item $\beta$，无量纲：吸水强度。
\end{itemize}

\subsection{植物的分类}

我们被要求研究植物群落中物种多样性与植物群落对干旱的响应之间的关系，因此必须在模型中包含若干种不同的植物。然而，求解多个植物种群随时间变化的模型是困难的，因为这些模型由若干相互耦合的微分方程以及大量参数构成。为在计算效率与模型精度之间取得平衡，我们根据植物应对干旱的不同策略将其划分为三个类别，并在每个类别中选取一到两种植物作为代表。

基于高晓燕 [4] 的研究，我们将植物划分为如下三类：

\begin{itemize}
\item \textbf{避旱型植物（Drought-Escape Plant）}：这类植物具有较高的固有生长率，这使它们能够在干旱发生之前完成其生命周期。作为快速生长的代价，这类植物不擅长应对干旱。一旦干旱来临，它们的地上部分会迅速枯萎以减少水分消耗。典型的避旱型植物为尼泊尔常春藤（\textit{Hedera nepalensis}）与小叶女贞（\textit{Ligustrum quihoui} Carr.）。
\item \textbf{御旱型植物（Drought-Avoiding Plant）}：这类植物具有中等的固有生长率，但它们通常具备抵御干旱的能力。当土壤湿度下降时，这类植物可以通过关闭气孔或其他方式来减少水分损失 [7]，从而能够在一段时间内于干旱中继续生长。典型的御旱型植物包括马尾松（\textit{Pinus massoniana} Lamb）与银白杨（\textit{Populus alba}）。
\item \textbf{耐旱型植物（Drought-Tolerant Plant）}：这类植物具有相当低的固有生长率。其中一些在整整一年之内仅生长几厘米。这类植物体内往往具有储水的特殊结构，使其能够从一年中为数不多的降雨中吸收大量水分，从而维持生长。我们选取漠蒿（\textit{Artemisia desertorum}）作为代表物种。
\end{itemize}

基于上述分析，我们发现植物的生长率与其抵御干旱的能力之间存在权衡关系。于是我们可以通过固有生长率与最适土壤湿度来区分这些植物。以固有生长率作为横轴、以最适土壤湿度作为纵轴，我们可以标出三类植物所处的位置，如图 4 所示。

\subsection{数据收集与参数估计}

上文我们将植物物种划分为三个类别，并选取了五种代表性植物。为模拟它们在不同土壤湿度下的生长，我们需要收集或估计它们的固有生长率以及最适于它们生长的土壤湿度。通过对 [5] 的研究，我们可以获得相关数据，如表 2 所示。

\begin{itemize}
\item 尼泊尔常春藤：最适土壤湿度 69\%，固有生长率 $4.57\times 10^{-3}\,\mathrm{day}^{-1}$；
\item 小叶女贞：最适土壤湿度 62\%，固有生长率 $3.79\times 10^{-3}\,\mathrm{day}^{-1}$；
\item 马尾松：最适土壤湿度 45\%，固有生长率 $9.33\times 10^{-4}\,\mathrm{day}^{-1}$；
\item 银白杨：最适土壤湿度 53\%，固有生长率 $1.15\times 10^{-3}\,\mathrm{day}^{-1}$；
\item 漠蒿：最适土壤湿度 25\%，固有生长率 $3.15\times 10^{-4}\,\mathrm{day}^{-1}$。
\end{itemize}

我们无法直接从文献中获得固有生长率，因此需要对其进行估计。下面我们以马尾松为例说明估计过程，并讨论结果的合理性。文献 [6] 指出，这株高 2.5 m、树干半径为 0.09 m 的树木在一个月内长高了 7 cm，而树干半径没有变化。考虑到木材的密度是一个常数，马尾松的固有生长率可用下式估计：

\begin{equation}
r_{\mathrm{pop.al.}}
=\frac{\Delta m}{m_{0}\times 30\ \mathrm{day}}
=\frac{\Delta V}{V_{0}\times 30\ \mathrm{day}}
=\frac{7}{250\times 30}
=9.33\times 10^{-4}\ \mathrm{day}^{-1}
\label{eq:growthrate}
\end{equation}

其中 $\Delta m$ 与 $\Delta V$ 分别为生物量增量与体积增量，$m_0$ 与 $V_0$ 分别为初始生物量与初始体积；由于木材密度为常数，故生物量之比可用体积之比代替，由此得到上式中的数值结果。

表 2 中这五种植物的固有生长率与最适土壤湿度，与我们在 3.2 节中所阐述的特征相一致：耐旱型植物生长最慢，避旱型植物生长最快。并且固有生长率越高，对干旱的耐受能力越低。这表明我们所估计的参数是合理的。

\section{任务一：GIID 模型}

在本节中，我们将建立 GIID 模型，以描述植物的生长、植物与环境之间的相互作用，以及不规则干旱的影响。

\subsection{植物的生长与竞争}

考虑到植物种群往往难以直接计数，因此选择植物生物量来建立植物生长方程更为合适。

首先，我们考虑理想情形，即植物在空间与其他资源不受限制的情况下可以无限生长。令 $r_i$ 表示第 $i$ 种植物的固有生长率，则其生物量 $x_i$ 应满足如下方程：

\begin{equation}
\frac{dx_i(t)}{dt}=x_i(t)r_i
\label{eq:expgrowth}
\end{equation}

由于生长环境是有限的，并且根据假设 3，所有植物共享一个共同的环境容量 $Env$，因此固有生长率将线性下降，即 $r_i(t)=r_i-x_i/Env$。于是我们引入 Logistic 生长模型，植物生长方程应改写为如下形式：

\begin{equation}
\frac{dx_i(t)}{dt}=x_i(t)r_i\left(1-\frac{x_i(t)}{Env}\right)
\label{eq:logistic}
\end{equation}

此外，我们应选取植物的初始生物量。我们假设各类植物的初始生物量相同，均为 $x_0$：

\begin{equation}
x_i(0)=x_0
\label{eq:init}
\end{equation}

最终，我们得到如下一阶非线性微分方程：

\begin{equation}
\begin{cases}
\dfrac{dx_i(t)}{dt}=x_i(t)r_i\left(1-\dfrac{x_i(t)}{Env}\right)\\[1.2em]
x_i(0)=x_0
\end{cases}
\label{eq:single}
\end{equation}

通过求解上述微分方程，我们可以得到植物生物量与时间之间的关系：

\begin{equation}
x_i(t)=\frac{Env}{1+\left(\dfrac{Env}{x_i(0)}-1\right)\exp(-r_it)}
\label{eq:solution}
\end{equation}

当多种植物生活在同一区域时，不同物种之间的相互作用不可避免。这是因为多种植物的生态位存在重叠，它们会竞争养分、阳光等自然资源。在这种情况下，一株植物的生长与繁殖便会受到其他植物的影响。

Lotka--Volterra 模型在处理物种之间的竞争问题时十分有用。其基本思想是：一个物种（第 $i$ 种植物）在生长过程中会受到来自所有其他物种的阻力。假设在一块土地上有 $n$ 种植物生长，则描述第 $i$ 种植物生物量的方程将改写为：

\begin{equation}
\frac{dx_i(t)}{dt}=x_i(t)r_i\left(1-\frac{x_i(t)}{Env}-\alpha\sum_{j=1,\,j\neq i}^{n}\frac{c_j}{c_i}x_j(t)\right)
\label{eq:lv}
\end{equation}

其中 $\alpha$ 为竞争强度；$c_i$ 为第 $i$ 种植物的竞争能力，它是一个归一化到 $[0,1]$ 范围内的常数。式中的 $c_j/c_i$ 项表示第 $j$ 种植物对第 $i$ 种植物所施加的影响。因此，植物的竞争能力越强，它在竞争中胜出的可能性就越大。

\subsection{植物与水之间的相互作用}

\subsubsection{水分与环境容量之间的关系}

上文我们设计了用于描述环境约束下植物生长以及物种之间相互作用的模型。然而，植物与更大范围环境之间的相互作用对竞争同样具有重要作用。根据假设 2，在干旱条件下，土壤湿度是决定植物能否存活的最重要因素。

设 $W$ 为土壤湿度，$bestW_i$ 表示第 $i$ 种植物所需的最适土壤湿度。我们用 $Env_i$ 表示受土壤湿度影响的第 $i$ 种植物的环境容量。

$Env_i$ 具有以下几个重要特征：

\begin{itemize}
\item 当 $W$ 等于 $bestW_i$ 时，$Env_i$ 等于 $Env$；
\item 当 $W$ 与 $bestW_i$ 之间存在巨大差异，即 $|W-bestW_i|\to+\infty$ 时，$Env_i$ 趋近于零。
\end{itemize}

根据中心极限定理，当植物物种数目趋于无穷时，每种植物的环境容量应服从正态分布，因此我们可以将环境容量表示为：

\begin{equation}
Env_i=Env\times\exp\left(-\frac{(W-bestW_i)^2}{MNW_i}\right)
\label{eq:envcap}
\end{equation}

其中 $MNW_i$ 是控制衰减速率的常数，在生态学中被称为第 $i$ 种植物的湿度生态位宽度。根据研究 [5]，当土壤湿度降至其最适值的 20\% 时，植物几乎无法存活。因此当 $W=0.2\,bestW_i$ 时，我们假设第 $i$ 种植物的生物量只能达到 $0.01\,Env$：

\begin{equation}
0.01\,Env=Env\times\exp\left(-\frac{(0.8\,bestW_i)^2}{MNW_i}\right)
\label{eq:mnw}
\end{equation}

求解上述方程，可得 $MNW_i=0.14\,bestW_i^{2}$。

\subsubsection{土壤湿度的变化}

在植物生长过程中，土壤湿度并非一个常数。由于植物吸收、降雨与蒸发的共同作用，土壤湿度 $W$ 会随时间变化。$W$ 的变化与每株植物所吸收的水分有关，同时也受到降水量 $p(t)$ 与蒸发速率的影响。因此，我们可以用如下微分方程来描述土壤湿度的变化：

\begin{equation}
\frac{dW(t)}{dt}=-\beta\sum_{i}x_i(t)W(t)-\gamma W(t)+p(t)
\label{eq:soilwater}
\end{equation}

其中：

\begin{itemize}
\item $\beta W(t)x_i(t)$ 是第 $i$ 种植物的吸水速率。
根据假设 1，每株植物吸收的水分与其生物量成正比。此外，相关研究 [7] 表明，当土壤湿度升高时，植物的吸水速率会增加。为简化模型，我们假设吸水速率也与土壤湿度 $W(t)$ 成正比。于是第 $i$ 种植物的吸水速率可写为 $\beta W(t)x_i(t)$，其中 $\beta$ 为吸水强度。
\item $\gamma W(t)$ 是蒸发速率。
文献 [7] 表明，在亚热带或温带地区，土壤的蒸发速率与土壤湿度成正比。因此蒸发速率可表示为 $\gamma W$，其中 $\gamma$ 是描述蒸发强度的常数。基于草地蒸发速率的监测数据，我们可以估计出 $\gamma=0.12\,\mathrm{day}^{-1}$。
\item $p(t)$ 表示降水量。
根据假设 3，这是为植物群落供水的唯一途径。它可以从气象记录中获得。
\end{itemize}

\subsection{不规则天气的模拟}

为了预测当植物群落遭遇各种不规则天气周期，以及在本应降水丰沛的时期遭遇干旱时，物种将如何变化，我们需要对天气状况进行模拟。因此，我们设计了一个函数来描述植物群落中周期性与非周期性降水，如下所示：

\begin{equation}
p(t)=\left[A_0+A_1\left|\sin\left(\frac{\pi}{T}t\right)\right|+A_2\,\delta(t-t_0)\right]\times\frac{1+I}{2}
\label{eq:precip}
\end{equation}

其中：

\begin{itemize}
\item $A_0$ 是全年最小降水量，同时也代表旱季的降水量。
\item $A_1\left|\sin\left(\frac{2\pi}{T}t\right)\right|$ 是一个周期为 $T$ 的函数，表示降水量的周期性变化。为使模拟更贴近现实，我们取 $T=1$ 年。
\item $A_2\,\delta(t-t_0)$ 表示随机发生的极端天气，$t_0$ 是极端天气发生的时刻。
\item $(1+I)/2$ 表示随机发生的干旱。在一年之内，$I$ 只能取 0 或 1。当 $I$ 取 0 时，当年的降水量为常年的一半，这意味着严重干旱的发生。我们假设每年发生干旱的概率为 30\%。
\end{itemize}

\subsection{模型的实现}

基于上文所述的关于植物生物量增长、竞争以及植物与更大环境之间相互作用的模型，我们最终可以得到模型的通用微分方程组：

\begin{equation}
\begin{cases}
\dfrac{dx_i(t)}{dt}=x_i(t)r_i\left(1-\dfrac{x_i(t)}{Env}-\alpha\sum\limits_{j=1,\,j\neq i}^{n}\dfrac{c_j}{c_i}x_j(t)\right)\\[1.4em]
\dfrac{dW(t)}{dt}=-\beta\sum\limits_{i}x_i(t)W(t)-\gamma W(t)+p(t)
\end{cases}
\label{eq:system}
\end{equation}

由于我们无法直接得到该方程组的解析解，因此需要使用一些数值方法。然而，该方程组并不适合使用 Matlab 中的 4--5 阶 Runge--Kutta 求解器，因为 $p(t)$ 项受到一些随机因素的影响。考虑到这只是一个一阶系统，我们通过将微分方程转化为差分方程对其进行离散化。接下来，我们以 $\Delta t=5$ 天为时间间隔对其进行递推求解。所得结果如下。

\subsubsection{单物种群落}

为观察各种植物在周期性天气波动与非周期性干旱下生物量随时间的变化，我们首先从御旱型中选取马尾松、从避旱型中选取尼泊尔常春藤。随后我们将它们单独种植，不考虑它们与其他植物的竞争。最终结果如下。

结合两幅图中的结果，我们可以识别出几个显著的现象：

\begin{itemize}
\item 两种植物的生物量都受到周期性天气变化的影响，从而产生波动。然而马尾松的波动幅度小于尼泊尔常春藤。这主要是因为马尾松是御旱型植物，对土壤湿度不太敏感。相反，尼泊尔常春藤作为避旱型植物，只有在较高的土壤湿度下才能生长。因此当干旱发生时，尼泊尔常春藤更容易受到影响。
\item 在干旱的恢复过程中，尼泊尔常春藤恢复得相对较快。因为它具有较大的固有生长率，这使得它在土壤湿度适宜其生长时能够迅速恢复生物量。
\end{itemize}

如上所述，这两种植物生物量随时间的变化存在明显差异。此外，它们与植物自身的最适土壤湿度与固有生长率相一致。这表明我们的模型是合理的。

\subsubsection{五物种群落}

为研究多植物群落如何应对干旱，我们模拟了五个物种相互作用时生物量随时间的变化。

由图 7，我们可以得出以下结论：

\begin{itemize}
\item 在我们所设定的条件下，两种御旱型植物（尼泊尔常春藤与小叶女贞）成为群落的优势种。而耐旱型植物（漠蒿）的生物量很小，只有在干旱出现、其他种群数量下降时才能生长。这主要是由于耐旱型植物的固有生长率极低，且其竞争能力较弱。
\item 属于同一类型的植物，其生物量随时间的变化相似；而属于不同类型植物的生物量变化则差异很大。
\item 在我们的模拟过程中，该植物群落经历了演替：起初，两种避旱型植物的生物量迅速增长并占据优势；然而随着御旱型植物（银白杨与马尾松）的生长，避旱型植物的生物量逐渐下降。这是因为避旱型植物的竞争能力较弱，且对土壤湿度的要求较高。
\item 在五植物群落中，总生物量的波动远小于单植物群落。
\end{itemize}
\section{任务二与任务三：物种多样性的作用}

通过上文对植物群落的建模，我们发现气候的周期性变化与不规则干旱都会导致植物生物量的变化。由上文中的图 6 与图 7 可以得出结论：五个物种共存时的生物量变动幅度远小于仅有一个物种存活时的变动幅度。因此，探究物种多样性的影响十分重要。

\subsection{任务二：物种数目的作用}

\subsubsection{定性分析}

为了定性分析物种数目的作用，我们分别模拟了单植物群落与五植物群落，随后观察这些群落中每个物种的生物量。为使我们的模拟能够更有效地描述群落的长期变化，我们将模拟时长设为 50 年。结果如下。

由图 8 与图 9，我们可以清楚地看到气候周期与不规则干旱的发生对不同群落的影响。

\begin{itemize}
\item 周期性气候与不规则干旱都会给生物量带来波动：周期性气候给各种植物的生物量带来小幅波动，而不规则干旱则带来更大的波动，因为它持续的时间更长。
\item 五植物群落的总生物量低于单植物群落，这是由于物种之间对水资源及其他资源的竞争。
\item 五个物种共存时总生物量的波动更小，这表明五植物群落比单植物群落具有更高的稳定性。
\end{itemize}

\subsubsection{定量分析}

为验证群落确实从植物的共存中获益，我们应对其特征进行定量描述。受一项关于中国西北地区群落研究的启发，我们将使用若干易于计算的代表性指标来评价群落。

对一个群落而言最重要的指标是总生物量，它反映了群落中的植物是否生长繁茂。若第 $i$ 种植物的生物量记为 $x_i$，则群落中的总生物量 $TB$ 为所有植物生物量之和。由于群落的总生物量随时间不断变化，我们取总生物量的均值作为群落繁茂程度的度量。

\begin{equation}
TB=\frac{1}{T_{\max}}\int_{0}^{T_{\max}}\sum_{i=1}^{n}x_i(t)\,dt
\label{eq:tb}
\end{equation}

其中 $T_{\max}$ 是运行模拟程序所需的总时长。

我们知道，群落总生物量随时间的变化与群落的稳定性有关。总生物量的变化越小，群落在遭遇不规则干旱时的生物量损失就越小，这意味着群落越稳定。一般而言，我们可以用总生物量的标准差来度量其变化，但标准差会受到数据均值与量纲（计量单位）的影响。为消除这一影响，我们定义了一个无量纲指标：总生物量的变异系数（CoV, Coefficient of Variation）。我们可以用总生物量的标准差除以均值来计算该指标：

\begin{equation}
CoV=\frac{std(TB)}{\overline{TB}}
=\frac{\sqrt{\int_{0}^{T_{\max}}\left(TB(t)-\overline{TB}\right)^{2}dt}}{\overline{TB}}
\label{eq:cov}
\end{equation}

在物理学中，我们通常用“熵”来度量系统的无序程度。Shannon--Wiener 指数 [9] 正是受熵的概念启发而提出的物种均匀度度量：在一个群落中，每个物种的平均数量为 $\bar{x}_i=\int_{0}^{T_{\max}}x_i(t)/T_{\max}\,dt$，而总平均生物量即上文已定义的 $TB$。于是第 $i$ 个物种所占的比例为 $\bar{x}_i/TB$。Shannon--Wiener 指数（S-W 指数）定义为：

\begin{equation}
H=\sum_{i=1}^{n}-\frac{\bar{x}_i}{TB}\log\frac{\bar{x}_i}{TB}
\label{eq:sw}
\end{equation}

根据其表达式可以发现，当每个物种所占比例相同时，即 $\bar{x}_i=TB/n$ 时，S-W 指数达到最大值。而各物种的分布越不均匀，S-W 指数就越小。

S-W 指数度量的是群落的绝对均匀度，由于其最大值为 $\log n$，故定义 Pielou 指数来度量群落的相对均匀度：

\begin{equation}
J=\frac{H}{\log n}
\label{eq:pielou}
\end{equation}

Pielou 指数是介于 $[0,1]$ 之间的数值，它度量的是群落当前的均匀度与群落所能达到的最大均匀度之间的关系。

\subsubsection{模型的实现与验证}

我们在 2.4 节中收集了五个植物物种的数据，其中两个为避旱型植物，两个为御旱型植物，一个为耐旱型植物。现在，我们将一至五个物种分别种植到群落中，以此改变群落中植物的数目。随后我们将模拟各种植物生物量随时间的变化，并利用结果计算上文所定义的各项指标。我们将尝试植物的所有可能组合，例如当物种数目 $n=2$ 时，共有 $C_5^2=10$ 种组合。

我们首先确定了总生物量与物种数目之间的关系，如图 10 所示。随后我们观察了群落稳定性（以生物量的 CoV 度量）与物种数目之间的关系，如图 11 所示。图中的每个点代表一个物种组合。直线连接的是群落中具有不同物种数目的各物种组合指标均值。

从这些图中我们可以注意到两个显著的现象：

\begin{itemize}
\item 随着物种数目的增加，总生物量的均值先增大后减小。
当物种数目为二时，总生物量最大。进一步地，当物种数目从一增加到二时，总生物量变得更高。这是由于某些物种之间的互补效应，其机制将在下文讨论。而当物种数目从二增加到五时，总生物量变得更低。这是因为群落中的物种过多，加强了对水资源及其他资源的竞争，而这种竞争压缩了每个植物物种的生存空间。
\item 随着物种数目的持续增加，总生物量的 CoV 呈下降趋势。
这表明物种更多的群落其总生物量的波动更小。也就是说，它们更稳定，也更抗干旱。
\end{itemize}

下图（图 12）展示了 S-W 指数与物种数目之间的关系，图 13 展示了 Pielou 指数与物种数目之间的关系。

从这些图中我们可以发现，S-W 指数与 Pielou 指数的取值随着物种数目的增加而上升。当物种数目达到 5 时，Pielou 指数达到 0.87。这表明群落中物种的分布已经非常接近均匀分布。因此我们得出结论：群落中物种数目越多，各物种的均衡发展越好。这是因为当物种数目较多时，种间竞争加剧，使得任何一种植物都难以拥有远超其他植物的生物量。

根据我们的结论，尽管更高的物种数目可能导致群落总生物量的下降，但总生物量的变化将会减小。此外，群落抗干旱的能力将会提高。与此同时，群落中物种的分布将更加均匀，每种植物都将拥有足够的生存与发展空间。这些结论与相关研究的结果相一致，因此它反映了我们模型的合理性。与单植物群落的情形相比，当存在两种植物时总生物量略有增加，而抗干旱能力显著提高，因此至少需要两种不同的植物才能使群落获益。

\subsection{任务三：物种类型的作用}

上文我们探究了物种数目与总生物量之间的关系，也描述了物种数目与稳定性之间的关系。然而从图 10 与图 11 同样可以明显看出，不同物种组合的总生物量与 CoV 存在显著差异。例如，若群落中存在两个不同物种，最低的总生物量为 200 kg，最高的总生物量为 400 kg。因此，物种的类型也能对我们的模型产生很大影响。为分析这一作用，我们选取了三种典型的物种组合，每种组合均由两个不同物种构成，并将对每种物种组合分别进行模拟。

\subsubsection{三种代表性物种组合}

如果群落中的两个物种具有相同的固有生长率、抗旱性与竞争能力，那么它们的生物量将以相同的模式变化。正是两个物种属性的差异导致了它们生物量趋势的不同。基于这一考虑，我们选取了三种代表性的物种组合：(a) 尼泊尔常春藤与小叶女贞，二者均属于避旱型植物，且在各种属性上差异很小；(b) 尼泊尔常春藤与银白杨，二者分别属于避旱型植物与御旱型植物，且它们的最适土壤湿度存在中等差异；(c) 尼泊尔常春藤与漠蒿，二者分别属于避旱型植物与耐旱型植物，且它们的最适土壤湿度差异很大。图 14 展示了这些物种组合各自生物量随时间的变化模式。

由图 14，我们可以总结出这三种不同物种组合的生长模式。

\begin{itemize}
\item 在由两种御旱型植物组合而成的群落 (a) 中，两种植物的生物量同步变化：二者在干旱来临时都急剧下降，并在干旱结束后迅速恢复。这是因为两种植物具有相近的最适土壤湿度，因此二者都显著受到干旱的影响。
\item 在由避旱型植物与御旱型植物组合而成的群落 (b) 中，两种植物的生物量同样显著受到干旱的影响。但两种植物的生物量呈现出相反的趋势。这是因为在降水充足的时期，土壤湿度较高，适宜避旱型植物生长，它可以成为优势种，而御旱型植物因受竞争限制而生长十分缓慢。在干旱时期，环境不适宜避旱型植物生长，其生物量显著下降，从而减轻了御旱型植物所面临的竞争压力，于是御旱型植物能够在干旱期间生长。
\item 在由避旱型植物与耐旱型植物组合而成的群落 (c) 中，几乎只有避旱型植物在生长，而耐旱型植物在该群落中近乎灭绝。这或许是因为两种植物生长所需的条件差异太大，以致它们在同一群落中共存十分困难。此外耐旱型植物的竞争能力较低，这导致了它的灭绝。
\end{itemize}

三种物种组合的总生物量随时间的变化趋势如图 15 所示。我们还分别计算了这三种组合的总生物量与生物量的 CoV，如表 3 所示。

\begin{itemize}
\item 组合 (a)：总生物量 206.5 kg，生物量的 CoV 为 0.1526；
\item 组合 (b)：总生物量 250.3 kg，生物量的 CoV 为 0.0632；
\item 组合 (c)：总生物量 220.8 kg，生物量的 CoV 为 0.1051。
\end{itemize}

\subsubsection{我们的结论}

由图 15 与表 3 我们可以发现，物种组合 (a) 与 (c) 的生物量波动相当大，它们在干旱发生时甚至损失了近一半的总生物量。相反，物种组合 (b) 的生物量波动较小，其总生物量也高于另外两种物种组合。这是因为物种组合 (b) 由避旱型植物与御旱型植物构成，它们的最适土壤湿度不同，从而使这两个物种的优势能够相互补充。在降水充足的时期与干旱时期，群落中具有不同的优势种，因而群落的总生物量相对稳定。相比之下，在物种组合 (c) 中，两种植物差异过大，无法共存。因此我们得出结论：物种组合对群落的总生物量与稳定性都有很大影响。总体而言，群落中各植物之间土壤湿度的差异越显著，植物群落越容易在剧烈变化的环境中存活，群落的稳定性与总生物量通常也越高。当然，各植物之间的差异也不应过大，因为差异过大会导致其中一个物种灭绝。

\section{任务四、任务五与任务六：环境变化对物种的影响}

在本节中，我们将讨论外部环境变化对群落所造成的影响，这些变化可以表现为干旱频率或强度的改变、生境侵占、环境污染等等。

\subsection{任务四：干旱强度与频率的变化}

现在，我们讨论干旱强度与频率的变化对群落的影响。在上文用于模拟降水的方程（式 11）中，我们将干旱概率设为 30\%，并将干旱年的降水量设为常年的 $1/2$。当干旱频率增加时，我们将干旱概率调整为 50\%；当干旱强度增加时，我们将干旱年的降水量改为常年的 $1/3$。在改变干旱的频率与强度之后，我们模拟了五植物群落中各物种生物量随时间的变化，并绘制了群落的总生物量曲线。

可以看到，遭受更严重干旱的群落与此前的群落存在很大差异。由图 16 与图 17，我们可以得出两个结论：

\begin{itemize}
\item 干旱强度与频率的增加引起群落中物种分布的变化。
御旱型植物与耐旱型植物的生物量均下降，而它们的下降为耐旱型植物的生存留出了空间。因此耐旱型植物的生物量在干旱期间显著增加。
\item 在极端干旱气候下，群落生物量的波动仍然很小（不超过 10\%）。
这归因于耐旱型植物在干旱期间的生长，它表明了五植物群落具有很强的稳定性。
\end{itemize}

接下来，我们应当定量刻画干旱强度与频率的增加究竟如何影响群落的总生物量与稳定性，以及物种多样性的增加是否真能帮助群落抵御干旱。于是我们改变群落中植物的数目，将一至五个物种分别放入群落。我们计算总生物量以度量群落的活力，计算总生物量的 CoV 以描述其稳定性。图 18 与图 19 给出了我们的计算结果。在这次模拟中我们同样尝试了所有物种组合，每个点代表一个物种组合，直线连接的是所有物种组合的均值。

由上列各图，我们可以发现：

\begin{itemize}
\item 随着干旱的加剧，群落的总生物量必然下降。
当干旱条件取初始值（干旱概率 30\%，干旱年降水减少至 $1/2$）时，群落的总生物量最大；当提高干旱频率时，群落的总生物量下降；当同时提高干旱频率与干旱强度时，群落的总生物量进一步下降。这是因为干旱环境不适宜大多数植物生长，并且可能加剧群落内部对水资源的竞争。
\item 干旱越严重，物种多样性对群落的贡献越大。
从总生物量曲线的趋势看，当干旱条件取初始值时，总生物量随物种数目的增加先增大后减小。当干旱频率提高时，总生物量仍然先增大后减小，但增大的幅度变大、减小的幅度变小。当干旱频率与干旱强度同时提高时，总生物量随物种数目的增加而增加。这可以用上文提到的互补效应来解释：在干旱发生期间，耐旱型植物能够生长，以补偿其他植物的生物量损失，而干旱越剧烈，这种互补效应就越显著。
\item 干旱加剧了群落中生物量的波动，但物种数目的增加总能减小这种波动。
从总生物量 CoV 的趋势看，干旱越剧烈，CoV 越高。这是因为严重干旱通常伴随降雨的急剧变化，这也会使群落的生物量迅速变化。无论干旱程度如何，CoV 总是随群落中物种数目的增加而下降。也就是说，物种数目的增加总是带来群落更高的稳定性。
\end{itemize}

此外，为模拟干旱频率的降低，我们将干旱发生的概率设为 10\%，干旱年的降水量仍设为常年的 50\%。我们绘制了群落总生物量随物种数目的变化如下。

此时，群落的生物量随物种数目的增加而持续下降，因为此时干旱发生得较少，多物种群落中的互补效应并不显著。相反，多个物种之间的竞争限制了总生物量的增加。

\subsection{任务五：生境缩减与环境污染的影响}

在本节中，我们将研究生境缩减与污染的影响。可供植物生存的土地总量因采伐而不断被挤压，进而导致群落环境容量的下降。此外，由于过度排放，一些土地受到污染，使植物难以生存。因此群落中可供植物利用的实际土地将会减少，这同样会导致环境容量的降低。于是，生境缩减与污染的后果都可以归因于环境容量的下降。

因此，我们将逐步降低植物群落的环境容量，并观察这一过程中植物生物量的变化。我们将环境容量设定为一个阶跃函数，以模拟渐进式的生境缩减或环境污染：

\begin{equation}
Env=
\begin{cases}
1000, & 0\le t<10\\
600, & 10\le t<20\\
500, & 20\le t<30\\
400, & 30\le t<40\\
300, & 40\le t\le 50
\end{cases}
\label{eq:stepenv}
\end{equation}

其中时间 $t$ 的单位为年。该函数表示：在 $[0,10]$ 年内，生境未缩减、环境未受污染；而在 $[10,20]$ 年内，环境容量降至原来的 60\%，以此类推。最终，在 $[40,50]$ 年内，环境容量降至原来的 30\%。我们模拟了单植物群落与五植物群落中各种植物生物量随时间的变化，并将后四个阶段的平均生物量与 $Env=1000$（环境未受污染且生境未缩减）时的平均生物量进行比较。我们计算了生物量的损失率，结果如图 21 与图 22 所示。

从这些图中，我们可以得出三个显著结论：

\begin{itemize}
\item 当环境容量下降时，生态系统的总生物量随之下降。
当环境容量降至原来的 70\% 时，单植物群落与五植物群落分别损失了 72.45\% 与 18.16\% 的生物量。
\item 物种多样性能够帮助群落抵御环境污染。
单植物群落易于受到环境污染或生境缩减的伤害。当环境容量降至原来的 40\% 时，其生物量已经损失了 60\%。当该群落遭遇干旱时，生物量维持在较低水平，群落几乎丧失了从干旱影响中恢复的能力。相反，当环境容量降至初始值的 30\% 时，五植物群落仅损失了 19\% 的生物量，并且在群落受到干旱影响之后，各物种的生物量仍能恢复。这表明多个物种的共存不仅使群落对干旱具有更好的抵抗力，而且对环境污染与生境缩减也具有更好的抵抗力。
\end{itemize}

\subsection{任务六：关于群落存续性的进一步讨论}

\subsubsection{确保群落的存续}

在上述分析中，我们得到了描述植物生物量变化的微分方程模型，其中考虑了干旱与不规则天气周期的影响。基于全部结论，我们发现植物群落的长期存续主要取决于三个层面：种群层面、群落层面与生态系统层面。

\begin{itemize}
\item \textbf{种群层面} 我们发现各种植物的抗旱能力各不相同。考虑到未来干旱不可避免地将会加剧，为确保植物群落的存续，我们最好选择最适土壤湿度相对较低的植物，例如耐旱型植物。在不同的土壤湿度下，优势种会从避旱型植物转变为耐旱型植物。因此，可以适度种植耐旱型植物，以提高干旱来临时群落的稳定性。
\item \textbf{群落层面} 我们发现多物种群落比单物种群落更稳定。因此，我们应确保一个植物群落中至少存在两种不同的植物。与此同时，我们发现不同植物的特征存在差异，因此在选择物种时应格外注意。我们应选择差异相对较大的植物，使它们能够为植物群落提供生态位互补，并使其在干旱等不利天气面前具有恢复力。然而，我们应避免使两个物种之间的差异过大，以防其中一种植物在竞争中被淘汰。
\item \textbf{生态系统层面} 我们发现污染与生境缩减会对植物群落产生较大的影响，因此我们应呼吁人们保护环境、减少污染排放，并对入侵物种等可能破坏植物生境的因素保持警惕。
\end{itemize}

\subsubsection{群落对更大范围环境的影响}

\begin{itemize}
\item 植物群落是生态系统的一部分，其中物种的数目与类型也会影响生态系统的稳定性。
\item 由于许多植物具有庞大的根系，当它们生长繁茂时可以起到固土作用，从而防治土壤沙漠化。
\item 植物是生态系统的生产者，因此一个稳定的植物群落能够为该区域提供必要的能量，从而维持各种生物在干旱条件下的生存。
\end{itemize}

因此，植物群落对环境的影响不容忽视。

\section{敏感性分析}

在上文中，有若干未知参数并非通过估计获得，而是由我们人为设定。为说明这些参数的变化如何影响模拟结果，我们通过改变这些参数的取值并观察模拟结果的变化来进行敏感性分析。

\subsection{植物吸水速率 $\beta$}

我们改变植物吸水速率 $\beta$ 的取值，并绘制群落中各物种生物量随时间的变化如下。

由图 23 可以看出：随着 $\beta$ 的增大，植物群落出现了两个显著变化：

\begin{itemize}
\item 群落中的总生物量持续下降。
在 $\beta$ 从 $10^{-3}$ 变化到 $2\times 10^{-2}$ 的过程中，总生物量下降了约 32\%。这是因为随着 $\beta$ 的增大，对水资源的竞争越激烈，环境中可供利用的水资源越少，因此植物的环境容量不断下降。
\item $\beta$ 越高，植物对干旱的响应越强烈。
这是因为植物的吸水速率越高，干旱期间的土壤湿度就越低。于是干旱对植物环境容量的影响就越大。
\end{itemize}

\subsection{植物竞争强度 $\alpha$}

竞争强度 $\alpha$ 是影响植物之间关系的重要因素。我们改变 $\alpha$ 的取值，并绘制群落中各物种生物量的变化曲线如下。

由图 24 可以看出，随着 $\alpha$ 的增大，植物群落同样出现了两个显著变化：

\begin{itemize}
\item $\alpha$ 的增大会导致生物量显著下降。
这是因为 $\alpha$ 越大，植物之间的竞争越强，植物之间会相互抑制，从而导致生物量的下降。
\item 随着 $\alpha$ 的增大，一些植物将无法在群落中存活。
当 $\alpha=10^{-3}$ 时，无法存活的是耐旱型植物；当 $\alpha=10^{-2}$ 时，御旱型植物也无法存活。这就是生物学中的竞争排斥原理。当竞争强度较小时，物种可以共存；而当竞争强度增大时，不适应环境的物种将逐渐被排除出当前群落，直到群落中只剩下一个或两个“最强”的物种。
\end{itemize}

\subsection{气候条件}

在上述所有模拟中，我们采用的是典型亚热带地区的气候条件，因此设定 $A_1=30$、$A_2=98$。事实上，降水参数的取值可能会显著影响群落中种群的分布。此处我们改变 $A_1$ 与 $A_2$ 的取值，并绘制群落中生物量随时间的变化如下：

\begin{itemize}
\item 图 25(a) 代表初始的亚热带气候；
\item 图 25(b) 代表典型的热带雨林气候，降水量更大且波动更小，其中避旱型植物因具有极高的生长率而占据优势；
\item 图 25(c) 代表典型的热带沙漠气候。该群落的总生物量显著降低，耐旱型植物占据优势，因为它们能够在土壤湿度很低的环境中存活。
\end{itemize}

综上所述，物种之间的竞争强度、植物的吸水速率以及气候模式都是影响群落中植物生长的关键因素。在将模型应用于实践时，必须对这些参数进行更准确的估计与测量。

\section{模型评价}

\subsection{优势}

\begin{itemize}
\item 我们详细描述了植物与环境之间相互作用的机制。
我们没有采用“植物生长率与降水量直接成正比”这一过于粗糙的假设，而是通过假设不同植物以不同速率从土壤中吸收水分、从而导致土壤湿度下降，并且土壤湿度又影响植物生长，来对相互作用进行建模。这使我们的模型更加贴近现实，也更具说服力。
\item 我们对植物群落的分析在定性与定量两个层面上展开。
我们不仅利用生物量曲线对植物生长过程进行定性分析，还计算了多种指标，例如植物群落的总生物量以及总生物量的 CoV，以定量度量植物群落的稳定性。这些分析使我们更容易解释群落为何能从多样性中获益。
\item 我们对实验结果的解释十分详尽。
对于我们开展的每一次模拟，我们不仅描述了结果，还努力追溯直观结果背后的原因。此外，我们运用一些生物学理论来解释我们的结论。
\end{itemize}

\subsection{不足与可能的改进}

\begin{itemize}
\item 我们对土壤湿度变化的描述过于简单。
我们假设植物的吸收速率与土壤湿度成正比，并且水的蒸发速率也与土壤湿度成正比，这些假设可能并不十分合理。在实践中，这些关系可能需要借助更多的实验或研究成果加以进一步讨论。
\item 我们使用自建的降水函数来模拟天气变化，而没有使用真实的气候数据。
我们自建的降水函数无法模拟现实天气变化的全部细节，这可能使我们的模型产生偏差。在模型的实际应用中，可以直接查找历史降水数据。
\item 我们忽略了除降水之外的所有因素。
在现实中，植物生长还可能受到温度、日照等其他因素的影响，如果我们希望进一步提高模型的精度，就需要仔细考虑这些因素的变化。
\end{itemize}

\section*{参考文献}

\begin{enumerate}
\item[{[1]}] H. J. W. MUTSAERS, A Dynamic Equation for Plant Interaction and Application to Yield-density-time Relations, Annals of Botany, Volume 64, Issue 5, November 1989, Pages 521--531.
\item[{[2]}] Ehleringer, J.R., Phillips, S.L., Schuster, W.S.F. et al. Differential utilization of summer rains by desert plants. Oecologia 88, 430--434 (1991).
\item[{[3]}] SARAH E. PARK, LAURENCE R. BENJAMIN, ANDREW R. WATKINSON, The Theory and Application of Plant Competition Models: an Agronomic Perspective, Annals of Botany, Volume 92, Issue 6, December 2003, Pages 741--748.
\item[{[4]}] Gao Xiaoyan. Study on drought resistance and its physiological mechanism in several species of Sedum[D]. Inner Mongolia Agricultural University, 2009.
\item[{[5]}] Zhao Xuelin, Wang Taitan, Meng Wenting et al. Soil moisture status and water quantity monitoring of Yang Chai and Black Sagebrush in Mao Wu Su sandy area[J]. Xinjiang Agricultural Science, 2021, 58(05): 937--946.
\item[{[6]}] Li XL, Liu YJ, Sun Fei et al. Spatial and temporal patterns of soil moisture and its response to rainfall variation in the clay layer of the interdune mound of white thorn scrub in the Ulan Buh Desert[J]. Soil and Water Conservation Bulletin, 2022, 42(06): 39--46.
\item[{[7]}] Kirschner GK, Xiao TT, Blilou I. Rooting in the Desert: A Developmental Overview on Desert Plants. Genes. 2021; 12(5): 709.
\item[{[8]}] Yang XG, Liu CHH, Wang L et al. Ecological restoration and reconstruction of desert grasslands: water-mediated system response, ecological thresholds and reciprocal feedbacks driven by artificial vegetation[J]. Journal of Ecology, 2023, 43(01): 95--104.
\item[{[9]}] Zhang JC, Zhang JX, Yuan HB et al. Plant community types and their diversity in Kumtag Desert[J]. Grassland Science, 2012, 29(10): 1581--1588.
\end{enumerate}

\end{document}
